% ---- pg-495 （书页 483） ----

种鸢尾属植物分开. 例如, $\pi_1$ 中的花瓣宽度大大小于 $\pi_3$ 中的花瓣宽度.

达罗赫和莫西曼(Darroch \& Mosimann, 见文献[6])认为鸢尾属植物可以用“形状”变量来刻画, 而不必借助花的大小. 令 $Y_1=X_1/X_2$ 为萼片形状, $Y_2=X_3/X_4$ 为花瓣形状, 将变量 $Y_1$ 和 $Y_2$ 用于判别将在练习 11.28 中加以探究.

\begin{mvfigure}
\begin{tikzpicture}[xscale=1.9,yscale=1.4]
% 外框
\draw (0.5,0) rectangle (3.75,2.68);
% 纵轴刻度
\foreach \vy in {0,0.5,1.0,1.5,2.0,2.5}{%
  \draw (0.5,\vy) -- (0.62,\vy);
  \node[left,font=\small] at (0.5,\vy) {\vy};}
% 横轴刻度标签
\node[below,font=\small] at (1,0) {$\pi_1$};
\node[below,font=\small] at (2,0) {$\pi_2$};
\node[below,font=\small] at (3,0) {$\pi_3$};
\node[below,font=\small] at (2.1,-0.33) {组};
\node[rotate=90,font=\small] at (0.02,1.34) {花瓣宽度};
% 第一组
\draw (0.81,0.19) rectangle (1.19,0.33);
\draw (0.81,0.27) -- (1.19,0.27);
\draw (1,0.43) -- (1,0.33);
\draw (0.9,0.43) -- (1.1,0.43);
\draw (1,0.19) -- (1,0.09);
\draw (0.9,0.09) -- (1.1,0.09);
\node[font=\small] at (1,0.47) {$*$};
\node[font=\small] at (1,0.62) {$*$};
% 第二组
\draw (1.81,1.18) rectangle (2.19,1.52);
\draw (1.81,1.33) -- (2.19,1.33);
\draw (2,1.85) -- (2,1.52);
\draw (1.9,1.85) -- (2.1,1.85);
\draw (2,1.18) -- (2,0.93);
\draw (1.9,0.93) -- (2.1,0.93);
% 第三组
\draw (2.81,1.82) rectangle (3.19,2.33);
\draw (2.81,2.03) -- (3.19,2.03);
\draw (3,2.50) -- (3,2.33);
\draw (2.9,2.50) -- (3.1,2.50);
\draw (3,1.82) -- (3,1.54);
\draw (2.9,1.54) -- (3.1,1.54);
\end{tikzpicture}
\phantomsection\label{fig:11.11}\mvcap{图 11.11\quad 三种鸢尾属植物的花瓣宽度的盒形图}
\end{mvfigure}

在判别分析中, 选用适当的变量往往很困难. 像本例中所给出的总结, 可以帮助研究者根据评估一种方法分类效果好坏的最终准则作出合理而简单的选择. \hfill$\blacksquare$

我们着重讨论线性判别法则(11-52)或(11-56), 许多商业计算机程序也都以这一法则为基础. 虽然线性判别法则有简单结构, 但必须记住, 它是根据多元正态性和相同协方差矩阵这两个很强的假定推出来的. 在将某个线性分类法则付诸应用之前, 必须对这些尝试性假定按照先是多元正态性, 然后是相同协方差的顺序进行检验. 如果这两个假定中有一个被破坏或两个都被破坏, 那么只要事先对数据做适当变换, 仍有可能得到改进的分类.

二次法则是线性判别函数分类的备选方案. 这种法则适用于正态性成立而相同协方差阵的假定被严重破坏的情形. 然而同线性法则相比, 正态性假定对二次法则来说似乎是更重要的条件. 如果对使用一次还是二次法则拿不准, 可以同时采用两种法则分类, 然后用拉亨布鲁克提留方法来考察它们的错误率.

\section{对多个总体进行判别的费希尔方法}  % 原书 11.6
\phantomsection\label{sec:11.6}

费希尔曾建议将他的判别方法(见 11.3 节)推广到多个总体的情况. 隐藏在费希尔判别分析背后的动机, 是需要获得某种合理的总体表示, 这种表示只包含观察值的少数几个线性组合, 诸如 $a_1'\mathbf{x}$, $a_2'\mathbf{x}$ 和 $a_3'\mathbf{x}$. 当人们对采用(1)目视检查法或(2)图描述法分离若干个总体感兴趣时, 这种方法具有下列优点:

1. 将数量很大的特征变量减少到相对说来较少几个线性组合, 从而形成维数较小的 $g$ 个总体的方便的表示. 当然, 最优分类所需的某些信息会丢失, 除非这些总体的均值完全位于所选的低维空间中.

2. 对前两个或三个线性组合(判别量)的均值作图, 将有助于展示总体之间的关系及可能的分组情况.

3. 作前两个判别量样本值的散布图, 有助于发现离群值和数据中其他异常情况.
